\chapter*{Acknowledgements}\label{chapter:ack}
\addcontentsline{toc}{chapter}{Acknowledgements}  
\markboth{}{} %% removes header from pages
% The acknowledgements is where you thank everyone who have contributed to the thesis. It is mandatory to thank your supervisor, co-supervisor, your department and other parties that you have cooperated with; fellow students, a company etc.  

% The acknowledgements is usually written in a more personal tone than other parts of the thesis, and you should use the pronoun \enquote{I} when you write. F.ex. I would like to thank Dag Langemyhr and Martin Helsø for their passion for \LaTeX, their dedication to elegance and detail in typography, and for reading the package documentation at \emph{The Comprehensive TeX Archive Network} (\hyperlink{https://www.ctan.org/}{CTAN})\footnote{Note that CTAN is the central place for all kinds of material around TeX, with 6649 packages as of October 2024.} so that the rest of us can worry less about layout and more about \hyperlink{https://en.wikipedia.org/wiki/Millennium_Prize_Problems}{The Millennium Prize Problems}\footnote{Footnotes are usually not used in mathematical texts, in particular not for citations! More on that later. In the meantime, use Oria via \hyperlink{https://ub.uio.no}{ub.uio.no} to find books about The Millennium Prize Problems}.

% In the acknowledgements it is custom to include important events that have happened during the period where you were working on the project. Sometimes, this can be done with humor and creativity. But be extremely cautious! Humor changes over time, your life and perspective changes, and you never know who will read this part of your thesis. What seems funny now might be extremely insulting or embarrassing in a couple of years\footnote{As a rule of thumb, try to refrain from being too creative and funny in your thesis. Write a blog, diary or send letters to a friend if you must.}.

% In fact, the acknowledgements is the part of the thesis that \emph{everyone} who picks up your thesis will read, not just the examiner or the interested researcher. I recommend that you visualize your audience when you feel most creative: your supervisor, your parents, your grandmother, your ex, your future boss, your future co-workers, your future kids, and, perhaps most importantly, your future students.

% The fact that your audience might be larger than you think, can not be stressed enough. And this applies to the entire thesis. Your thesis will eventually be deposited in and made openly available to the entire world via the institutional repository \hyperlink{https://duo.uio.no}{DUO}. Because of the UiO-address, it gets a high rating on search engines and will be very easy to find for interested researchers all over the world. Moreover, most researchers subscribe to alerts when new citations of their articles or books are published, which means that most of the researchers that you have cited are notified about your thesis, and they might read it, push it to their students, find errors, notice sloppy work and referencing, pick up interesting results and more.

I would like to thank my supervisor Jørgen Vold Rennemo for all the help, support and guidance during this project. The supervising meetings have been a joy, and I have learned so much in the past year.

I would like to thank the entire 11th floor of the Mathematics Department. This is a special place filled with great people. I have learned so much from fellow students, the PhD candidates, and the professors, and the time spent here has been so much fun thanks to all of you. A special thanks to Aleksander Leraand, who went from being an inspiration as a senior master student, to being the biggest aid in my thesis apart from my supervisor. Your help has been invaluable.

I want to thank Sander Oppen for proofreading this thesis and fixing some of my more embarrassing grammar mistakes.

I would also like to thank my family and friends for all the support and encouragement during my studies. I would especially like to thank my wonderful girlfriend Elin for her love, patience, and understanding throughout this journey. This thesis would not have been possible without her.

% Finally, I would like to thank the Department of Mathematics at the University of Oslo for providing a great learning environment and resources for my studies.